% ============================================================
% PÁGINA 1 — CAPA
% ============================================================
\pagenumbering{gobble}
\thispagestyle{coverpage}
\setlength{\parindent}{0pt}

\vspace*{2cm}

\begin{center}
\textbf{\MakeUppercase{%
Monitoramento e Diagnóstico de Eventos Disruptivos nos Portos
Brasileiros: Detecção de Anomalias Operacionais e Análise de
Exposição Comercial}}
\end{center}

\vspace{2cm}

\hfill
\begin{minipage}{0.55\textwidth}
{\setstretch{1.0}%
Trabalho de Conclusão de Curso apresentado como parte dos requisitos
para obtenção do grau de Especialista em Ciência de Dados e
Inteligência Artificial.

\vspace{1cm}

\textbf{Aluno:} Fábio Kouri Paim\\
\textbf{Orientador:} Prof. Alex Lopes Pereira
\par}
\end{minipage}

\vfill

\begin{center}
Brasília -- DF\\
Abril/2026
\end{center}

\newpage

% ============================================================
% PÁGINA 2 — FOLHA DE ROSTO / RESUMO
% ============================================================
\thispagestyle{coverpage}

\begin{center}
\textbf{Monitoramento e Diagnóstico de Eventos Disruptivos nos Portos Brasileiros:
Detecção de Anomalias Operacionais e Análise de Exposição Comercial}

\vspace{0.8cm}

Fábio Kouri Paim\\
\textit{Escola Nacional de Administração Pública (ENAP)}\\
\textit{MBA em Ciência de Dados}
\end{center}

\vspace{1cm}

{\setstretch{1.0}%
\noindent\textbf{Resumo.}
Este artigo propõe um sistema integrado de detecção e classificação de anomalias operacionais em portos brasileiros, articulando três bases públicas: o Anuário Estatístico Aquaviário (ANTAQ), o PortWatch/FMI e o ComexStat/MDIC. Esse cruzamento liga a atividade portuária no nível da atracação individual aos fluxos de exportação detalhados por produto, unidade da federação de origem e país de destino, conectando a operação física à cadeia produtiva exportadora. A partir de séries temporais semanais de 18 portos, modelos XGBoost com correção AR(1) aprendem o padrão operacional normal de cada porto, e um ensemble de três métodos estatísticos identifica desvios significativos. Um sistema de classificação dual categoriza cada anomalia quanto à sua abrangência (isolada, nacional ou global), distinguindo choques domésticos de choques no comércio marítimo global. Validado contra eventos históricos conhecidos, o sistema identifica 1.609 anomalias no período 2015--2026. A integração com dados de exportação por cadeia produtiva mapeia pares porto-cadeia de fragilidade estrutural e estados com dependência elevada de um único porto. Os resultados sugerem que a articulação de dados públicos já disponíveis pode oferecer a gestores públicos uma base analítica reprodutível para avaliar a exposição de cadeias exportadoras a disrupções. Como entrega prática, o trabalho disponibiliza um painel interativo de monitoramento com atualização semanal.
\par}

\vspace{0.5cm}

\begin{flushright}
\textbf{Palavras-chave:} detecção de anomalias; portos brasileiros; machine learning;\\
cadeia de suprimentos; vulnerabilidade comercial.
\end{flushright}

\newpage
\setlength{\parindent}{1.25cm}
\pagenumbering{arabic}
\setcounter{page}{1}
